\documentclass[12pt]{article}
\usepackage[utf8]{inputenc}
\usepackage[spanish]{babel}
\usepackage[a4paper, margin=2cm]{geometry}
\usepackage{tikz}
\usetikzlibrary{babel}
\usepackage[most]{tcolorbox}
\usepackage{amsmath}

% Usar fuente sin serifa para un aspecto más moderno (tipo infografía)
\renewcommand{\familydefault}{\sfdefault}

% Definición de colores con temática médica/científica
\definecolor{medblue}{RGB}{43, 108, 176}
\definecolor{medteal}{RGB}{49, 151, 149}
\definecolor{medlight}{RGB}{235, 248, 255}
\definecolor{meddark}{RGB}{26, 32, 44}
\definecolor{ecggreen}{RGB}{16, 185, 129} % Verde estilo monitor de signos vitales

\begin{document}
\pagestyle{empty}

\begin{tcolorbox}[
    enhanced, % Habilita características avanzadas como las sombras
    colback=medlight!30, 
    colframe=medblue, 
    title={\Large \textbf{Sistemas de Coordenadas: Electrodos de ECG}}, 
    halign title=center, 
    fonttitle=\bfseries, 
    arc=4mm, 
    boxrule=1.5pt,
    shadow={2mm}{-2mm}{0mm}{black!30!white}
]

    \vspace{0.2cm}
    \begin{tcolorbox}[colback=white, colframe=medteal, arc=2mm, boxrule=1.2pt, title=\textbf{Caso Clínico / Problema}]
        Se colocan dos electrodos en el pecho de un paciente. En un sistema de referencia cartesiano trazado sobre la piel, las posiciones de los electrodos están en \textbf{(3, 5) cm} y \textbf{(9, 13) cm}. Calcule la distancia recta entre ambos.
    \end{tcolorbox}

    \vspace{0.4cm}
    
    \begin{center}
    \begin{tikzpicture}[scale=0.55]
        % Cuadrícula de fondo
        \draw[step=1cm,gray!40,very thin] (-0.5,-0.5) grid (10.5,14.5);
        
        % Ejes cartesianos
        \draw[thick,->,meddark] (-0.5,0) -- (10.5,0) node[anchor=north west] {$x$ (cm)};
        \draw[thick,->,meddark] (0,-0.5) -- (0,14.5) node[anchor=south east] {$y$ (cm)};
        
        % Marcas numéricas
        \foreach \x in {2,4,6,8,10} \draw (\x cm,0.2cm) -- (\x cm,-0.2cm) node[anchor=south, yshift=0.1cm] {\footnotesize $\x$};
        \foreach \y in {2,4,6,8,10,12,14} \draw (0.2cm,\y cm) -- (-0.2cm,\y cm) node[anchor=east] {\footnotesize $\y$};
        \node[anchor=north east] at (0,0) {\footnotesize $0$};
        
        % Coordenadas de los puntos
        \coordinate (E1) at (3,5);
        \coordinate (E2) at (9,13);
        \coordinate (Aux) at (9,5); % Punto auxiliar para el triángulo rectángulo
        
        % Triángulo rectángulo (componentes)
        \draw[dashed, medteal, thick] (E1) -- (Aux) node[midway, below, fill=medlight!30, inner sep=2pt] {$\Delta x = 9 - 3 = 6$};
        \draw[dashed, medteal, thick] (Aux) -- (E2) node[midway, right, fill=medlight!30, inner sep=2pt] {$\Delta y = 13 - 5 = 8$};
        
        % Distancia de Electrodos (Línea principal)
        \draw[line width=3pt, ecggreen, cap=round] (E1) -- (E2) node[midway, above left, text=ecggreen!80!black, font=\bfseries] {$d = ?$ (Distancia)};
        
        % Puntos destacados (Electrodos)
        \filldraw[meddark] (E1) circle (4pt) node[anchor=north west, font=\bfseries, fill=white, inner sep=1pt] {Electrodo 1 (3,5)};
        \filldraw[meddark] (E2) circle (4pt) node[anchor=south east, font=\bfseries, fill=white, inner sep=1pt] {Electrodo 2 (9,13)};
        
    \end{tikzpicture}
    \end{center}

    \vspace{0.4cm}
    \begin{tcolorbox}[colback=medteal!10, colframe=medteal, arc=2mm, boxrule=1.2pt, title=\textbf{Desarrollo Matemático y Solución}]
        Para encontrar la distancia lineal entre los dos electrodos de ECG, aplicamos la fórmula de distancia:
        \begin{align*}
            d &= \sqrt{(x_2 - x_1)^2 + (y_2 - y_1)^2} \\[1ex]
            d &= \sqrt{(9 - 3)^2 + (13 - 5)^2} = \sqrt{(6)^2 + (8)^2} \\[1ex]
            d &= \sqrt{36 + 64} = \sqrt{100} = \mathbf{10 \text{ cm}}
        \end{align*}
        \textbf{Resultado:} La distancia recta entre los dos electrodos sobre el pecho del paciente es de \textbf{10 cm}.
    \end{tcolorbox}
    
\end{tcolorbox}

\end{document}